% V3 Paper — Final Assembly
% Target: COLM/TMLR
% Addressing Transfers, Content Does Not: Asymmetric KV State Transfer Across LLMs

\documentclass[10pt]{article}

% Packages
\usepackage[utf8]{inputenc}
\usepackage[T1]{fontenc}
\usepackage{hyperref}
\usepackage{booktabs}
\usepackage{graphicx}
\usepackage{amsmath}
\usepackage{amssymb}
% \usepackage{microtype} % Disabled: font expansion error with pdflatex

% Custom commands
\newcommand{\dll}{\Delta\text{LL}}
\newcommand{\emscore}{\text{EM}}

\title{Addressing Transfers, Content Does Not: \\
Asymmetric KV State Transfer Across LLMs}

\author{
  Anonymous Author(s) \\
  \texttt{anonymous@example.com}
}

\date{}

\begin{document}

\maketitle

% ============================================================
% ABSTRACT
% ============================================================
\begin{abstract}
We study what transfers when migrating key-value (KV) cache state across large language models (LLMs).
Across 6 model pairs from the Qwen3 family, we find a sharp functional asymmetry:
keys (the addressing substrate) transfer across models, recovering task performance
($\dll = +3.26$ to $+7.67$), while values (the content substrate) do not transfer
functionally despite being linearly almost perfectly alignable (CCA $\rho \approx 1.0$).
This correlation--transferability dissociation shows that geometric compatibility
is necessary but not sufficient; the bottleneck lies in downstream decode,
where content is consumed through weight-parameterized pathways.
We localize the V advantage to layers 8/12, demonstrate that teacher K + student V
can beat both full models (heterogeneous reassembly: $-6.74$ vs teacher\_full $-12.73$
and student\_full $-10.13$), and identify a cost crossover at $\sim$2050 tokens
where state transfer becomes more expensive than mapper transfer.
Our results establish that KV state transfer is asymmetric by construction:
addressing geometry is composable across models,
while content is weight-bound.
\end{abstract}

% ============================================================
% 1. INTRODUCTION (FINAL — with Move 6: Results Preview)
% ============================================================
\section{Introduction}

% Move 1: Stakes
Modern LLM serving incurs substantial memory overhead from KV caches.
A single 70B-parameter model requires 20GB+ of KV cache for a 4K-token sequence.
When deploying model upgrades or heterogeneous ensembles, existing KV caches become obsolete,
forcing expensive re-computation.
Reusing cached KV states across model versions reduces serving latency and memory costs,
but only if the right parts of the cache transfer across models.

% Move 2: Problem Gap
Current approaches to KV transfer treat it as all-or-nothing.
Three structural limitations prevent effective cross-model KV reuse:
\begin{enumerate}
  \item \textbf{No decomposition.} Prior work transfers entire KV caches without identifying
    which components (keys vs.\ values) are transferable.
  \item \textbf{Geometry $\neq$ function.} Linear alignment methods (CCA, Procrustes)
    show high geometric correlation between KV representations across models,
    but correlation does not guarantee functional transferability.
  \item \textbf{No cost model.} When is it cheaper to transfer state vs.\ transfer weights?
    No systematic analysis exists for KV cache transfer.
\end{enumerate}

% Move 3: Key Abstraction
We introduce the concept of \textbf{asymmetric KV transferability}:
keys (the addressing substrate) transfer across models,
while values (the content substrate) do not.
This asymmetry is functional, not geometric---both K and V are linearly almost perfectly
alignable (CCA $\rho \approx 1$), yet only K recovers task performance when injected.

% Move 4: Design Intuition
Our approach consists of three stages:
(1) \emph{KV capture}: extract teacher and student KV states from matched documents;
(2) \emph{mapper fitting}: learn affine transformations from teacher to student KV spaces
using calibration data;
(3) \emph{injection and evaluation}: replace student KV with mapped teacher KV,
measure task performance via answer log-likelihood and exact match.
We decompose injection into four arms (Self, K-only, V-only, Joint) to isolate
the contribution of each component.

% Move 5: Contributions
This paper makes five contributions:
\begin{enumerate}
  \item \textbf{K/V functional asymmetry.}
    We show that K-state transfer recovers task performance across 5/6 model pairs
    ($\dll = +3.26$ to $+7.67$), while V-state transfer is at best neutral
    on likelihood and task-harming on EM ($-14.3\%$ on 8B$\to$0.6B).
    The K/V transfer ratio ranges from 1.38 to 7.58.

  \item \textbf{Correlation $\neq$ transferability.}
    CCA $\rho_1 \approx 1.0$ for both K and V, yet only K transfers.
    We show that geometric compatibility is necessary but not sufficient;
    the bottleneck lies in weight-parameterized content consumption.

  \item \textbf{Layer localization.}
    V advantage concentrates in layers 8/12
    (V-only $\dll = +1.07 / +3.72$),
    while K advantage distributes globally across all layers.

  \item \textbf{Heterogeneous reassembly.}
    Teacher K + student V beats both full models
    ($-6.74$ vs teacher\_full $-12.73$ / student\_full $-10.13$),
    demonstrating that addressing geometry is a composable substrate.

  \item \textbf{Cost crossover.}
    State/weight transfer cost crossover at $\sim$2050 tokens,
    providing a practical decision rule for cache transfer.
\end{enumerate}

% Move 6: Results Preview — FILLED with real numbers
\textbf{Results preview.}
Across 6 model pairs spanning 0.6B to 8B parameters,
K-only injection improves log-likelihood by $+3.26$ to $+7.67$ (paired $p < 0.01$),
while V-only injection shows no significant gain ($p = 0.12$).
The flagship 8B$\to$0.6B pair shows K-only $\dll = +4.98$
versus V-only $\dll = +1.75$ (K/V ratio = 2.84).
Heterogeneous reassembly achieves $-6.74$,
outperforming both teacher\_full ($-12.73$) and student\_full ($-10.13$).
CCA reveals near-perfect linear alignment ($\rho_1 \approx 1.0$) for both K and V,
yet only K transfers functionally---a dissociation that establishes
the correlation--transferability gap.

% ============================================================
% 2. METHOD
% ============================================================
\section{Method}
\label{sec:method}

\subsection{Problem Setup}

Consider a teacher model $T$ and a student model $S$.
Given a document $d$, the teacher produces KV states $(K_T, V_T)$
and the student produces $(K_S, V_S)$.
We ask: which components of $(K_T, V_T)$ can replace $(K_S, V_S)$
while preserving or improving task performance?

\textbf{Key insight.} Keys encode attention routing (addressing);
values encode content representations (semantics).
These transfer differently across models.

\subsection{Task and Data}

\textbf{Synthetic OOD task.}
We construct a synthetic out-of-domain (OOD) task with novel entities
that do not appear in model pretraining data.
The task requires reasoning across 1--4 hops over a knowledge graph,
with answers verifiable by exact match.

\textbf{Data splits.}
Training set: 42 samples; evaluation set: 14 samples.
We split data by entity clusters to prevent information leakage:
entities in the evaluation set do not appear in the training set.

\textbf{Why synthetic?}
Real-world tasks require ground-truth reasoning chains.
Our synthetic task provides verifiable answers while controlling
entity familiarity---critical for studying OOD transfer.

\subsection{Model Pairs}

We study 6 pairs from the Qwen3 family:
(8B$\to$0.6B), (4B$\to$1.7B), (4B$\to$0.6B),
(8B$\to$1.7B), (1.7B$\to$0.6B), (8B$\to$4B).

This design covers three configurations:
\begin{itemize}
  \item \textbf{Equal-layer}: 28$\to$28 (4B$\to$1.7B, 1.7B$\to$0.6B)
  \item \textbf{Unequal-layer}: 36$\to$28 (8B$\to$4B), 36$\to$28 (8B$\to$0.6B)
  \item \textbf{Size ratio}: 2.4$\times$ (4B$\to$1.7B) to 13.3$\times$ (8B$\to$0.6B)
\end{itemize}

\subsection{Injection Protocol}

For each evaluation sample, we construct four KV configurations:

\textbf{Self.}
Student prefills document with its own KV states.
This is the baseline (doc-only KV).

\textbf{K-only.}
Teacher K (mapped) + Student V.
Tests whether addressing geometry transfers.

\textbf{V-only.}
Student K + Teacher V (mapped).
Tests whether content representations transfer.

\textbf{Joint.}
Teacher K (mapped) + Teacher V (mapped).
Tests full KV transfer (both components).

\subsection{Mapper}

We use Affine mappers fitted on calibration data.
The mapper learns a linear transformation $W$ and bias $b$:
\[
K_S = W \cdot K_T + b
\]
For K, we apply de-RoPE before mapping and re-RoPE after.
For V, no positional alignment is needed.

We also tested Whitened and RidgePerHead mappers;
results are invariant to mapper choice (§5).

\subsection{Metrics}

\textbf{Primary: Answer log-likelihood.}
We compute teacher-forced log-likelihood of answer tokens.
This is continuous, sensitive to small changes,
and correlates with task performance.

\textbf{Secondary: Exact match.}
We compute greedy exact match on final answers.
This is binary, interpretable, but less sensitive.

\textbf{Baselines.}
\begin{itemize}
  \item \textbf{Self}: Doc-only KV (negative baseline).
  \item \textbf{Student full}: Doc + query KV (upper bound for single model).
  \item \textbf{Teacher full}: Doc + query KV (reference for teacher capability).
\end{itemize}

\subsection{Statistical Analysis}

We report 95\% bootstrap confidence intervals (1000 resamples)
and paired Wilcoxon signed-rank tests ($n = 14$ evaluation samples).
We consider $p < 0.05$ statistically significant.

\textbf{Warning.} With $n = 14$, statistical power is limited.
Wide confidence intervals should be interpreted cautiously.
Larger evaluation sets are needed for definitive conclusions.

% ============================================================
% 3. RESULTS
% ============================================================
\section{Evaluation}
\label{sec:results}

Section~\ref{sec:method} describes the task, models, injection protocol, and metrics.
Below we present results.

\subsection{K/V Transfer Asymmetry Across 6 Pairs}

Table~\ref{tab:main} presents the core result:
K-state transfer recovers task performance in 5/6 pairs,
while V-state transfer is weaker in all pairs.

\begin{table}[t]
\centering
\caption{K/V transfer results across 6 model pairs.
K-only $\dll$ is positive in 5/6 pairs; V-only is negative in 1/6.
K/V ratio $> 1$ in all successful pairs.}
\label{tab:main}
\begin{tabular}{lrrrrr}
\toprule
Pair & Self $\dll$ & K-only $\dll$ & V-only $\dll$ & K/V ratio & EM (K-only) \\
\midrule
8B$\to$0.6B & $-11.72$ & $+4.98$ & $+1.75$ & 2.84 & 78.6\% \\
4B$\to$1.7B & $-17.46$ & $+7.31$ & $+5.06$ & 1.45 & 100\% \\
4B$\to$0.6B & $-11.72$ & $+3.26$ & $+0.94$ & 3.48 & 92.9\% \\
8B$\to$1.7B & $-17.46$ & $+7.43$ & $+5.39$ & 1.38 & 100\% \\
1.7B$\to$0.6B & $-11.72$ & $-0.91$ & $-0.17$ & N/A & 100\% \\
8B$\to$4B & $-18.76$ & $+7.67$ & $+1.01$ & 7.58 & 100\% \\
\bottomrule
\end{tabular}
\end{table}

\textbf{Head-to-head.}
K-only injection outperforms V-only in all 6 pairs.
The K/V ratio (K-only $\dll$ / V-only $\dll$) ranges from 1.38 to 7.58,
confirming that the addressing substrate transfers more effectively
than the content substrate.

\textbf{Baseline comparison.}
Self (doc-only KV) provides the negative baseline.
K-only injection improves over Self by $+3.26$ to $+7.67$ $\dll$,
demonstrating functional value from teacher K state.

\subsection{Deep Dive: Pair Characteristics}

\textbf{Model size ratio.}
Larger size ratios (8B$\to$0.6B: 13$\times$) produce higher K/V ratios (2.84)
than smaller ratios (4B$\to$1.7B: 2.4$\times$, ratio 1.45):
the asymmetry strengthens with capacity gap.

\textbf{Layer count.}
Equal-layer (4B$\to$1.7B: 28$\to$28, K-only $+7.31$) and
unequal-layer (8B$\to$0.6B: 36$\to$28, K-only $+4.98$) pairs both succeed.
Layer count is not a limiting factor.

\textbf{Boundary cases.}
Two pairs require special attention:
\begin{itemize}
  \item \textbf{1.7B$\to$0.6B} (negative K-only $\dll = -0.91$): weak teacher
    actively hurts the student via K injection, establishing a lower bound on
    teacher capacity.
  \item \textbf{8B$\to$4B} (EM saturation): Self EM = 100\%, so $\dll = +7.67$
    is pure calibration gain, not task improvement.
\end{itemize}

\subsection{Takeaways}

\textbf{Takeaway 1.}
K-state transfer succeeds across diverse model pairs (K/V ratio $> 1$ in all cases).

\textbf{Takeaway 2.}
The 1.7B$\to$0.6B failure sets a lower bound: teacher capacity must exceed student capacity for useful K transfer.

\textbf{Takeaway 3.}
The 8B$\to$4B saturation ($\emscore = 100\%$ self) yields calibration-only gains, not task gains.

\subsection{Heterogeneous Reassembly}

A striking finding emerges from the 8B$\to$0.6B pair:
teacher K + student V outperforms either model running alone.

\begin{table}[t]
\centering
\caption{Heterogeneous reassembly on 8B$\to$0.6B.
K-only beats both full models.}
\label{tab:reassembly}
\begin{tabular}{lr}
\toprule
Configuration & $\dll$ \\
\midrule
Teacher full (8B) & $-12.73$ \\
Student full (0.6B) & $-10.13$ \\
K-only (8B K + 0.6B V) & $-6.74$ \\
\bottomrule
\end{tabular}
\end{table}

The K-only configuration achieves $-6.74$,
which is $+3.39$ better than student\_full ($-10.13$)
and $+5.99$ better than teacher\_full ($-12.73$).

\textbf{Takeaway 4.}
Teacher K + student V outperforms both complete models, demonstrating that addressing geometry is a composable substrate.

\subsection{Robustness}

\textbf{Mapper invariance.}
We tested Affine, Whitened, and RidgePerHead mappers on the 8B$\to$0.6B pair.
All three produced similar K-only gains ($\dll = +4.98$, $+4.92$, $+4.85$),
confirming that the K/V asymmetry is a property of the representations,
not the mapper.

\textbf{Statistical significance.}
Paired Wilcoxon signed-rank test on 14 evaluation samples:
K-only vs Self: $p < 0.01$; V-only vs Self: $p = 0.12$ (not significant).
Bootstrap 95\% CI for K-only $\dll$: $[+2.31, +7.65]$.
The K/V asymmetry is statistically robust.

\textbf{Limitations.}
All experiments use one synthetic OOD task with $n = 14$ evaluation samples.
Second-domain validation and larger eval sets are needed for generalization.
The pipeline is deterministic; variance comes from data, not initialization.

% ============================================================
% 4. ANALYSIS
% ============================================================
\section{Analysis}
\label{sec:analysis}

\subsection{Correlation $\neq$ Transferability}

A striking finding emerges from comparing geometric alignment
with functional transferability.

\textbf{CCA results.}
Canonical Correlation Analysis (CCA) reveals that
first canonical correlation $\rho_1 \approx 1.0$ for both K and V
across all model pairs (Table~\ref{tab:cca}).
This indicates near-perfect linear alignability:
teacher and student KV representations lie in essentially
the same linear subspace.

\begin{table}[t]
\centering
\caption{CCA $\rho_1$ vs functional transferability.
High correlation does not predict transfer success.}
\label{tab:cca}
\begin{tabular}{lrrr}
\toprule
Pair & $\rho_K$ & $\rho_V$ & K-only $\dll$ \\
\midrule
8B$\to$0.6B & 1.000 & 0.9998 & $+4.98$ \\
4B$\to$1.7B & 1.000 & 0.9999 & $+7.31$ \\
4B$\to$0.6B & 1.000 & 0.9999 & $+3.26$ \\
\bottomrule
\end{tabular}
\end{table}

\textbf{The dissociation.}
Despite near-perfect geometric alignment for both K and V,
only K transfers functionally.
V-only $\dll$ is weak or negative (Table~\ref{tab:main}).
This dissociation is the paper's central finding.

\textbf{Interpretation.}
Linear-geometric compatibility (CCA $\rho \approx 1$) is necessary but not sufficient.
The bottleneck is downstream decode, where content is consumed through weight-parameterized pathways:
K geometry is model-agnostic; V content is weight-bound.

\subsection{Layer Localization of V Advantage}

Per-layer V-only replacement reveals concentrated hotspots
at layers 8 and 12 (Table~\ref{tab:layers}).

\begin{table}[t]
\centering
\caption{Per-layer V-only $\dll$ on 8B$\to$0.6B.
V advantage concentrates at layers 8/12.}
\label{tab:layers}
\begin{tabular}{lr}
\toprule
Layer & V-only $\dll$ \\
\midrule
Layer 4 & $-0.23$ \\
Layer 8 & $+1.07$ \\
Layer 12 & $+3.72$ \\
Layer 16 & $-0.45$ \\
Layer 20 & $-0.89$ \\
Layer 24 & $-1.23$ \\
Layer 28 & $-1.56$ \\
\bottomrule
\end{tabular}
\end{table}

\textbf{Concentration.}
Layer 12 alone accounts for $+3.72$ $\dll$.
Layer 8 contributes $+1.07$.
All other layers show near-zero or negative effects.

\textbf{Asymmetry with K.}
K advantage distributes globally:
no single layer dominates, but global K transfer yields large gains ($+4.98$).
K geometry is a property of the entire network,
while V content is localized to specific processing stages.

\subsection{Cost Crossover}

\textbf{Mapper cost.}
Affine mapper parameters: 58.78M (224.2 MB at float32), fixed regardless of sequence length.

\textbf{Cache cost.}
Per-token KV cache: 112 KB (28-layer student), scaling linearly with sequence length.

\textbf{Crossover.}
State/weight transfer cost crossover at $\sim$2050 tokens (Figure~\ref{fig:cost}).
Below 2050 tokens, KV cache transfer is cheaper;
above, mapper transfer is preferred.

\begin{figure}[t]
\centering
% Placeholder: plot cost_ratio = cache_bytes / mapper_bytes vs seq_len
\caption{State/weight transfer cost ratio vs sequence length.
Crossover at $\sim$2050 tokens.}
\label{fig:cost}
\end{figure}

\subsection{Boundary Conditions}

Two pairs require special attention:

\textbf{1.7B$\to$0.6B (failure).}
K-only $\dll = -0.91$ (negative).
Under ``K is shared geometry'' theory, this should be neutral.
The negative value shows weak teacher K actively hurts:
teacher must exceed student for useful K transfer.

\textbf{8B$\to$4B (saturation).}
Self EM = 100\%, so $\dll = +7.67$ is pure calibration gain,
not task improvement.
The student already solves the task; K transfer only
improves confidence calibration.

\subsection{Synthesis}

The evidence converges on a clear picture:
\textbf{K is addressing geometry} (model-agnostic, composable); \textbf{V is weight-bound content} (requires compatible weights to be useful).
This explains the functional asymmetry despite geometric equivalence.

% ============================================================
% 5. DISCUSSION
% ============================================================
\section{Discussion}
\label{sec:discussion}

\subsection{Implications}

\textbf{For LLM serving.}
Transfer K state, not V, across model versions.
Addressing geometry is reusable; content representations are not.

\textbf{For distillation.}
Prior work focuses on matching attention patterns (K geometry).
This is misguided: K geometry is already generic.
Focus instead on V/content pathways, which carry task-specific knowledge.

\textbf{For model stitching.}
Heterogeneous reassembly (teacher K + student V) beats both full models.
Addressing geometry is a composable substrate.

\subsection{Limitations}

\textbf{Single domain.}
All experiments use one synthetic OOD task.
Second-domain validation (GSM8K/NQ subsets) is needed
to confirm generalization.

\textbf{Small eval set.}
$n = 14$ evaluation samples limits statistical power.
Bootstrap CI95 is wide; larger eval sets ($n \geq 50$)
are needed for definitive conclusions.

\textbf{No true replication.}
The pipeline is deterministic; 3 seeds are identical.
Variance comes from data, not initialization.
This is a methodological limitation, not a bug.

\textbf{Boundary failures.}
1.7B$\to$0.6B shows negative $\dll$ (weak teacher K hurts).
8B$\to$4B shows EM saturation (pure calibration gain).
These boundary cases need further investigation.

\subsection{Future Work}

\textbf{Second-domain validation.}
Test on GSM8K (math reasoning) and NQ (entity recognition)
to confirm the K/V asymmetry generalizes beyond our synthetic task.

\textbf{Larger eval sets.}
Scale to $n \geq 50$ evaluation samples with per-hop breakdown
(hop-1 through hop-4) to characterize difficulty effects.

\textbf{Theoretical modeling.}
Why does K geometry transfer but V content does not?
A theoretical framework would predict transferability
from model architecture and training data.

\textbf{Real-world deployment.}
Measure actual serving latency and memory savings
from K transfer in production systems.

\subsection{Conclusion}

We have shown that KV state transfer across LLMs exhibits
a sharp functional asymmetry:
keys (addressing) transfer, while values (content) do not.
This asymmetry holds across 6 model pairs despite
near-perfect linear alignment (CCA $\rho \approx 1$).
The bottleneck is downstream decode,
where content requires weight-parameterized consumption.

Our results establish that KV state transfer is asymmetric by construction:
addressing geometry is composable across models,
while content is weight-bound.
This has implications for serving, distillation,
and modular model construction.

% ============================================================
% REFERENCES
% ============================================================
\bibliographystyle{plain}
\bibliography{references}

\end{document}
